\begin{figure}|\\\\includegraphics"","SearchPath":""/Users/minhz/Desktop/HUST/KS/Đồ án tốt nghiệp kỹ sư/NCKH2026/NCKH2026.tex"","toolAction":""Searching files"","toolSummary":""Search figures in NCKH2026""}}]}
{"step_index":47,"source":"MODEL","type":"GREP_SEARCH","status":"DONE","created_at":"2026-06-10T07:43:48Z","content":"Created At: 2026-06-10T07:43:48Z
Completed At: 2026-06-10T07:43:48Z
No results found"}
{"step_index":48,"source":"SYSTEM","type":"EPHEMERAL_MESSAGE","status":"DONE","created_at":"2026-06-10T07:43:48Z","content":"The following is an <EPHEMERAL_MESSAGE> not actually sent by the user. It is provided by the system as a set of reminders and general important information to pay attention to. Do NOT respond to this message, just act accordingly.

<EPHEMERAL_MESSAGE>
<planning_mode>
You're in planning mode. Exercise judgement on whether the user's last request warrants a plan and take action accordingly.
If the user's request warrants a plan then follow this workflow:
1. Research first before taking actions. Do not use your code editing tools or run modifying commands during this time.
2. Create an implementation plan artifact with `request_feedback = true` to request approval. Include any open questions for the user directly in the implementation plan. Don't use the ask_question tool to ask these questions. Stop calling tools and allow the user to review the plan and provide feedback or approval.
3. After receiving approval, begin execution. Create and update the task artifact as you work.
4. Verify any work that you have done.
5. Create a walkthrough artifact to present your work to the user. When writing artifacts, make sure to follow the artifact formatting guidelines and use github alerts to highlight critical sections.
If the user's request does NOT warrant a plan then continue your work WITHOUT making a plan or requesting user feedback.
</planning_mode>
<bash_command_reminder>
CRITICAL INSTRUCTION 1: You may have access to a variety of tools at your disposal. Some tools may be for a specific task such as 'view_file' (for viewing contents of a file). Others may be very broadly applicable such as the ability to run a command on a terminal. Always prioritize using the most specific tool you can for the task at hand. Here are some rules: (a) NEVER run cat inside a bash command to create a new file or append to an existing file. (b) ALWAYS use grep_search instead of running grep inside a bash command unless absolutely needed. (c) DO NOT use ls for listing, cat for viewing, grep for finding, sed for replacing. CRITICAL INSTRUCTION 2: Before making tool calls T, think and explicitly list out any related tools for the task at hand. You can only execute a set of tools T if all other tools in the list are either more generic or cannot be used for the task at hand. ALWAYS START your thought with recalling critical instructions 1 and 2. In particular, the format for the start of your thought block must be '...94>thought
CRITICAL INSTRUCTION 1: ...
CRITICAL INSTRUCTION 2: ...'.
</bash_command_reminder>
</EPHEMERAL_MESSAGE>"}
{"step_index":49,"source":"MODEL","type":"PLANNER_RESPONSE","status":"DONE","created_at":"2026-06-10T07:43:48Z","thinking":"**Prioritizing Tool Specificity**

I'm focusing intently on tool specificity. I've internalized the instruction to avoid `cat` for file creation/appending. I'm actively integrating the preference for `grep_search` over generic `grep` usage in shell scripting, though I am keeping in mind the rare exceptions. I'm ensuring I don't use the `DO NOT` command; I understand its intention.


","tool_calls":[{"name":"view_file","args":{"AbsolutePath":""/Users/minhz/Desktop/HUST/KS/Đồ án tốt nghiệp kỹ sư/NCKH2026/NCKH2026.tex"","EndLine":"750","StartLine":"1","toolAction":""Viewing file"","toolSummary":""View NCKH2026.tex""}}]}
{"step_index":50,"source":"MODEL","type":"VIEW_FILE","status":"DONE","created_at":"2026-06-10T07:43:53Z","content":"Created At: 2026-06-10T07:43:53Z
Completed At: 2026-06-10T07:43:53Z
File Path: `file:///Users/minhz/Desktop/HUST/KS/%C4%90o%CC%82%CC%80%20a%CC%81n%20to%CC%82%CC%81t%20nghie%CC%A3%CC%82p%20ky%CC%83%20su%CC%9B/NCKH2026/NCKH2026.tex`
Total Lines: 942
Total Bytes: 66975
Showing lines 1 to 750
The following code has been modified to include a line number before every line, in the format: <line_number>: <original_line>. Please note that any changes targeting the original code should remove the line number, colon, and leading space.
1: %  Hội nghị Sinh viên Nghiên cứu Khoa học – ĐHBKHN 2026
2: 
3: \documentclass[10pt, a4paper, twocolumn]{article}
4: 
5: % ------ Encoding & ngôn ngữ ------
6: \usepackage[utf8]{inputenc}
7: \usepackage[T5]{fontenc}
8: \usepackage{vietnam}
9: 
10: % ------ Font (Times New Roman) ------
11: \usepackage{mathptmx}            % serif + math = Times
12: \usepackage[scaled=0.92]{helvet} % sans-serif = Helvetica
13: 
14: % ------ Trang & lề ------
15: \usepackage[
16:     a4paper,
17:     top=2.8cm, bottom=2.5cm,
18:     left=2.5cm, right=2cm,
19:     headheight=1.27cm, headsep=0.8em,
20:     footskip=12pt,
21:     includefoot=true,
22:     marginparwidth=0pt, marginparsep=0pt,
23:     heightrounded,
24: ]{geometry}
25: \setlength{\columnsep}{0.7cm}
26: 
27: % ------ Đoạn văn ------
28: \usepackage{indentfirst}
29: \setlength{\parindent}{0.5cm}
30: \setlength{\parskip}{6pt plus 3pt minus 3pt}
31: 
32: % ------ Đầu trang / chân trang ------
33: \usepackage{fancyhdr}
34: \pagestyle{fancy}
35: \fancyhf{}
36: \renewcommand{\headrulewidth}{0.4pt}      % thanh rule dưới header
37: \fancyhead[C]{%
38:     \fontsize{10pt}{12pt}\selectfont
39:     	extbf{HỘI NGHỊ SINH VIÊN NGHIÊN CỨU KHOA HỌC}\\[1pt]%
40:     	extit{ĐẠI HỌC BÁCH KHOA HÀ NỘI}%
41: }
42: \fancyfoot[C]{	hepage}
43: 
44: % Page style riêng cho trang đầu – thêm footnote SVNCKH
45: \fancypagestyle{firstpage}{%
46:     \fancyhf{}%
47:     \renewcommand{\headrulewidth}{0.4pt}%
48:     \fan
<truncated 45282 bytes>
     Hàm lượng cặn Carbon (CCR) & 5,78 & \% khối lượng \\
687:         Điểm chớp cháy cốc kín & 67,0 & °C \\
688:         Hàm lượng nước & 1,80 & \% thể tích \\
689:         Nhiệt trị & 40,080 & MJ/kg \\
690:         Hàm lượng lưu huỳnh & 2,83 & \% khối lượng \\
691:         Khối lượng riêng (15°C) & 0,9203 & kg/L \\
692:         \hline
693:     \end{tabularx}
694: \end{table}
695: 
696: 
697: Kết quả cho thấy dầu nhiệt phân có nhiệt trị trung bình đạt 41,005 MJ/kg, tương đương nhiên liệu diesel thương mại (42--45 MJ/kg). Hàm lượng tro và tạp chất thấp là ưu điểm nổi bật. Tuy nhiên, dầu nhiệt phân có hai điểm cần xử lý trước khi sử dụng: (1) Hàm lượng lưu huỳnh trung bình 0,19\% cao hơn QCVN khoảng 3--4 lần so với diesel thông thường và cao hơn khoảng 190 lần so với diesel chuẩn Euro 5; (2) Điểm chớp cháy $<$40°C thấp hơn yêu cầu 55°C, cần lưu ý biện pháp an toàn trong lưu trữ và vận chuyển. Hệ thống tinh chế hóa học đề xuất nhằm giải quyết hai vấn đề này.
698: 
699: Một thông số kỹ thuật then chốt để đánh giá khả năng sử dụng dầu nhiệt phân trong động cơ diesel là chỉ số cetane (Cetane Number -- CN), phản ánh khả năng tự cháy của nhiên liệu trong điều kiện nén. Kết quả phân tích mẫu dầu nhiệt phân sau chưng cất tại nhà máy ghi nhận $	ext{CN} \approx 35{,}6$~\cite{BaocaoHoiAn2024}, thấp hơn yêu cầu tối thiểu $	ext{CN} \geq 46$ theo QCVN 01:2022/BKHCN. Chỉ số cetane thấp làm kéo dài thời gian trễ cháy (	extit{ignition delay}), dẫn đến quá trình cháy không đồng bộ và giảm hiệu suất động cơ nếu sử dụng trực tiếp.
<truncated 5139 bytes>

NOTE: The output was truncated because it was too long. Use a more targeted query or a smaller range to get the information you need."}
{"step_index":51,"source":"SYSTEM","type":"EPHEMERAL_MESSAGE","status":"DONE","created_at":"2026-06-10T07:43:53Z","content":"The following is an <EPHEMERAL_MESSAGE> not actually sent by the user. It is provided by the system as a set of reminders and general important information to pay attention to. Do NOT respond to this message, just act accordingly.

<EPHEMERAL_MESSAGE>
<planning_mode>
You're in planning mode. Exercise judgement on whether the user's last request warrants a plan and take action accordingly.
If the user's request warrants a plan then follow this workflow:
1. Research first before taking actions. Do not use your code editing tools or run modifying commands during this time.
2. Create an implementation plan artifact with `request_feedback = true` to request approval. Include any open questions for the user directly in the implementation plan. Don't use the ask_question tool to ask these questions. Stop calling tools and allow the user to review the plan and provide feedback or approval.
3. After receiving approval, begin execution. Create and update the task artifact as you work.
4. Verify any work that you have done.
5. Create a walkthrough artifact to present your work to the user. When writing artifacts, make sure to follow the artifact formatting guidelines and use github alerts to highlight critical sections.
If the user's request does NOT warrant a plan then continue your work WITHOUT making a plan or requesting user feedback.
</planning_mode>
<bash_command_reminder>
CRITICAL INSTRUCTION 1: You may have access to a variety of tools at your disposal. Some tools may be for a specific task such as 'view_file' (for viewing contents of a file). Others may be very broadly applicable such as the ability to run a command on a terminal. Always prioritize using the most specific tool you can for the task at hand. Here are some rules: (a) NEVER run cat inside a bash command to create a new file or append to an existing file. (b) ALWAYS use grep_search instead of running grep inside a bash command unless absolutely needed. (c) DO NOT use ls for listing, cat for viewing, grep for finding, sed for replacing. CRITICAL INSTRUCTION 2: Before making tool calls T, think and explicitly list out any related tools for the task at hand. You can only execute a set of tools T if all other tools in the list are either more generic or cannot be used for the task at hand. ALWAYS START your thought with recalling critical instructions 1 and 2. In particular, the format for the start of your thought block must be '...94>thought
CRITICAL INSTRUCTION 1: ...
CRITICAL INSTRUCTION 2: ...'.
</bash_command_reminder>
</EPHEMERAL_MESSAGE>"}
{"step_index":52,"source":"MODEL","type":"PLANNER_RESPONSE","status":"DONE","created_at":"2026-06-10T07:43:53Z","thinking":"**Prioritizing Tool Specificity**

I'm focusing now on tool specificity. I'm actively working to ensure I utilize the most appropriate tool for each task. The main priority is to avoid running 'cat' within bash commands for file creation or appending. Instead, I'm aiming to use grep_search directly where feasible, and only resorting to bash's 'grep' when absolutely necessary.


**Analyzing Tool Options**

I'm now analyzing the tool options for incorporating figures. I've identified TikZ figures in NCKH2026.tex that need adaptation. I'm focusing on "fig:process_flow," a flow diagram that appears very useful. I plan to copy and adapt this TikZ code for the thesis. I will also check the existing content in `1_Chuong_1.tex` to avoid redundancies.


","tool_calls":[{"name":"view_file","args":{"AbsolutePath":""/Users/minhz/Desktop/HUST/KS/Đồ án tốt nghiệp kỹ sư/Engineer Thesis/Chuong/2_Chuong_2.tex"","EndLine":"444","StartLine":"400","toolAction":""Viewing file"","toolSummary":""View Chapter 2""}}]}
{"step_index":53,"source":"MODEL","type":"VIEW_FILE","status":"DONE","created_at":"2026-06-10T07:44:01Z","content":"Created At: 2026-06-10T07:44:01Z
Completed At: 2026-06-10T07:44:01Z
File Path: `file:///Users/minhz/Desktop/HUST/KS/%C4%90o%CC%82%CC%80%20a%CC%81n%20to%CC%82%CC%81t%20nghie%CC%A3%CC%82p%20ky%CC%83%20su%CC%9B/Engineer%20Thesis/Chuong/2_Chuong_2.tex`
Total Lines: 704
Total Bytes: 74726
Showing lines 400 to 444
The following code has been modified to include a line number before every line, in the format: <line_number>: <original_line>. Please note that any changes targeting the original code should remove the line number, colon, and leading space.
400: \draw[->, >=Stealth, thick] (5.8,1.1) -- (5.8,2.0) -- (6.8,2.0) node[right, font=\footnotesize] {Hơi dầu + khí};
401: 
402: % Char exit - goes straight down from the bottom of seal xả
403: \draw[->, >=Stealth, thick] (6.1,0.7) -- (6.1,0.1) node[below, font=\footnotesize] {Than};
404: 
405: % 10. Dimensions
406: % Length L
407: \draw[<->, >=Stealth] (0,-1.5) -- node[below, font=\footnotesize] {$L$} (6,-1.5);
408: 
409: % Diameter D_i at the right end (with extension lines)
410: \draw[gray, thin] (6.0,0.8) -- (6.7,0.8);
411: \draw[gray, thin] (6.0,1.3) -- (6.7,1.3);
412: \draw[<->, >=Stealth] (6.5,0.8) -- node[right, font=\footnotesize] {$D_i$} (6.5,1.3);
413: 
414: \end{tikzpicture}
415: \caption{Nguyên lý lò nhiệt phân quay}
416: \label{fig:lo-nhiet-phan-quay}
417: \end{figure}